\documentclass[11pt]{article}
\usepackage[margin=1in]{geometry}
\usepackage{amsmath,amssymb,amsthm,mathtools}
\usepackage{hyperref}
\usepackage{enumitem}
\usepackage{booktabs}

\title{Verifiable Resource Certificates for Public Reversible Arithmetic Blocks in Quantum ECDLP}
\author{A. Rahman}
\date{\today}

\begin{document}
\maketitle

\begin{abstract}
Quantum resource estimates for the elliptic-curve discrete logarithm problem
(ECDLP) now shape cryptographic migration planning, blockchain security analysis,
and fault-tolerant architecture design. Recent work has moved in two complementary
directions: Babbush et al. give improved secp256k1 resource estimates supported by
zero-knowledge attestation while withholding sensitive circuit details, whereas Luo
et al. publish an explicit reversible modular-inversion construction based on the
extended Euclidean algorithm, reducing the logical-qubit footprint of prime-field
ECDLP and identifying gate count, depth, and architecture-aware implementation as
natural optimization targets. This note proposes a third disclosure model:
verifiable resource certificates for public reversible arithmetic blocks. A
certificate records a circuit commitment, gate basis, resource counts, input-output
specification, deterministic test generation, correctness transcript, and optional
proof artifact. We specialize the framework to modular inversion blocks
\(|x\rangle|0\rangle\mapsto |x\rangle|x^{-1}\bmod p\rangle\), for
\(x\in\mathbb F_p^\times\), which are central to affine-coordinate quantum ECDLP
implementations. We prove a basic soundness bound for hash-derived randomized
testing and outline a prototype verifier. The goal is not a new quantum attack,
but reproducible, comparable, and independently auditable quantum-ECDLP arithmetic
claims.
\end{abstract}

\tableofcontents

\section{Introduction}

TODO: expand this section.

\section{Background on Quantum ECDLP Arithmetic}

TODO: expand this section.

\section{Certificate Objects for Reversible Arithmetic Blocks}

TODO: expand this section.

\section{Deterministic Test Generation and Soundness}

TODO: expand this section.

\section{Certificate Specialization: Modular Inversion over Prime Fields}

TODO: expand this section.

\section{Prototype Verifier}

TODO: expand this section.

\section{Comparison with Existing Disclosure Models}

TODO: expand this section.

\section{Limitations and Responsible Scope}

TODO: expand this section.

\section{Future Work}

TODO: expand this section.


\bibliographystyle{alpha}
\bibliography{refs}
\end{document}
